Treating the UF as its own currency in the standardised measures changes the income measure far more than the value measure, and for different reasons. On the committed book the single-currency treatment of the same flows gives an income measure \repRatioLo{} to \repRatioHi{} times smaller and a value measure \repEveShortLo{} to \repEveShortHi{} percent smaller in the three variants with a net UF position of 30 percent of level 1 capital; over 30 books the income ratios are $\mainRatioNiiLng\pm\mainRatioNiiLngCi$, $\mainRatioNiiSht\pm\mainRatioNiiShtCi$ and $\mainRatioNiiOff\pm\mainRatioNiiOffCi$ and the value shortfalls $\mainEveShortLng\pm\mainEveShortLngCi$, $\mainEveShortSht\pm\mainEveShortShtCi$ and $\mainEveShortOff\pm\mainEveShortOffCi$ percent. The exact Shapley decomposition attributes the income gap to the indexation charge, which is the constant $\kappa|\PNUR|$, and to the pair rule, in a proportion of about two to one, and the value gap to the own curve and the own sizes, which nearly cancel, with the pair rule small except where the two legs move in opposite directions.

Two of the predictions and one further statement were refuted in whole or in part, and the refutations are results. The pair rule does not dominate the value measure. The Annex 1 income measure does not always exceed the Basel positive-part aggregation: it falls below it when the charge is smaller than the offset reduction, as Equation~\eqref{eq:charge-condition} states. A binding UF floor can raise a measure, in \floorsRaise{} of \floorsTotal{} comparisons after the curve is translated down, because it reduces the gain that the pair rule offsets. The bound that places the pair rule between the sum of the legs and the positive parts held in every one of the \hOneChecked{} comparisons, and the two classifications differ only at net UF positions of \hFourLo{} to \hFourHi{} percent of level 1 capital, through the value test.

The next steps are to run the decomposition on the flows of a real bank where one is available, to replace the translation of the curve by historical periods of negative real rates when their yields are in the raw data, to obtain a second implementation of Annex 1 to reconcile with, and to extend the five switches to the option treatment and to a third currency. The numbers are those of a synthetic bank over the curves of one day.
